% --- Archon physics formalization source begin ---
% archon:physics
% archon:covers IPhO2026Problems/problem_IPhO_2026_4_C_7.lean
% archon:source-report reports/ipho_2026/problem_IPhO_2026_4_C_7.source.json
% archon:problem-id IPhO_2026_4
% archon:part-id C.7
% archon:previous-part IPhO_2026_4_C_6
% archon:previous-part-policy natural_language_prerequisite_only

\chapter{Physics problem IPhO\_2026\_4\_C\_7}
\label{ch:IPhO2026Problems_problem_IPhO_2026_4_C_7}

\paragraph{Problem source.}
Water in the inner and outer cylinders exchanges heat radially through an
acrylic cylindrical wall.  Record T\_IC and T\_OC versus time.  The heat-flow
model is dQ/dt = (T\_OC - T\_IC)/R\_Th.  For radial Fourier conduction,
dQ/dt = -lambda*A*dT/dr.  Ignore apparatus heat capacity where instructed and
use the dimensions in Figure 17.

Current subquestion:
Combine the heat-flow relation and radial Fourier law to determine acrylic conductivity lambda.

\paragraph{Current subquestion.}
Combine the heat-flow relation and radial Fourier law to determine acrylic conductivity lambda.

\paragraph{Recorded answer/context.}
lambda = ln(r\_2/r\_1)/(2*pi*h*R\_Th). Official sample: lambda = 0.25 +/- 0.01 W/(m*K).

\paragraph{Figure/image path.}
/root/proposal\_for\_physic/science-mango/ipho\_2026\_source/image/E1\_page-14.png

\paragraph{Reusable previous-part conclusions.}
\begin{itemize}
\item Source C.6. Question: Determine the effective wall thermal resistance R\_Th from the C5 graph. Reusable conclusions: R\_Th = 1/(c\_0*m*slope). Official sample: R\_Th = 1.17 +/- 0.03 K/W. Policy: natural\_language\_prerequisite\_only; do\_not\_import\_Lean\_output
\end{itemize}

\paragraph{Formalization target.}
create a compiling Lean file with sorry bodies at `IPhO2026Problems/problem\_IPhO\_2026\_4\_C\_7.lean`. The Lean declarations must preserve the physical quantities, dimensions or dimensional roles, figure labels, governing-law hypotheses, and final relation expressed by this problem.
Use Mathlib/Physlib names found through LeanExplore where available. If a physics API is missing, introduce faithful local abstractions rather than scalar placeholder aliases.

\begin{definition}[powerDimension]
  \label{decl:physics:IPhO_2026_4_C_7:powerDimension}
  \lean{IPhO2026Problems.IPhO2026_4_C_7.powerDimension}
  The physical dimension of energy per unit time.
\end{definition}
\begin{proof}
  This is the defining declaration; unfolding it gives the stated typed data or relation.
\end{proof}

\begin{definition}[thermalResistanceDimension]
  \label{decl:physics:IPhO_2026_4_C_7:thermalResistanceDimension}
  \lean{IPhO2026Problems.IPhO2026_4_C_7.thermalResistanceDimension}
  \uses{decl:physics:IPhO_2026_4_C_7:powerDimension}
  The physical dimension kelvin per watt.
\end{definition}
\begin{proof}
  This is the defining declaration; unfolding it gives the stated typed data or relation.
\end{proof}

\begin{definition}[thermalConductivityDimension]
  \label{decl:physics:IPhO_2026_4_C_7:thermalConductivityDimension}
  \lean{IPhO2026Problems.IPhO2026_4_C_7.thermalConductivityDimension}
  \uses{decl:physics:IPhO_2026_4_C_7:powerDimension}
  The physical dimension watt per metre-kelvin.
\end{definition}
\begin{proof}
  This is the defining declaration; unfolding it gives the stated typed data or relation.
\end{proof}

\begin{definition}[specificHeatCapacityDimension]
  \label{decl:physics:IPhO_2026_4_C_7:specificHeatCapacityDimension}
  \lean{IPhO2026Problems.IPhO2026_4_C_7.specificHeatCapacityDimension}
  The physical dimension joule per kilogram-kelvin.
\end{definition}
\begin{proof}
  This is the defining declaration; unfolding it gives the stated typed data or relation.
\end{proof}

\begin{definition}[DimLength]
  \label{decl:physics:IPhO_2026_4_C_7:DimLength}
  \lean{IPhO2026Problems.IPhO2026_4_C_7.DimLength}
  A unit-independent length quantity.
\end{definition}
\begin{proof}
  This is the defining declaration; unfolding it gives the stated typed data or relation.
\end{proof}

\begin{definition}[DimMass]
  \label{decl:physics:IPhO_2026_4_C_7:DimMass}
  \lean{IPhO2026Problems.IPhO2026_4_C_7.DimMass}
  A unit-independent mass quantity.
\end{definition}
\begin{proof}
  This is the defining declaration; unfolding it gives the stated typed data or relation.
\end{proof}

\begin{definition}[DimTemperature]
  \label{decl:physics:IPhO_2026_4_C_7:DimTemperature}
  \lean{IPhO2026Problems.IPhO2026_4_C_7.DimTemperature}
  A unit-independent temperature quantity.
\end{definition}
\begin{proof}
  This is the defining declaration; unfolding it gives the stated typed data or relation.
\end{proof}

\begin{definition}[DimPower]
  \label{decl:physics:IPhO_2026_4_C_7:DimPower}
  \lean{IPhO2026Problems.IPhO2026_4_C_7.DimPower}
  \uses{decl:physics:IPhO_2026_4_C_7:powerDimension}
  A unit-independent signed heat-flow rate (power).
\end{definition}
\begin{proof}
  This is the defining declaration; unfolding it gives the stated typed data or relation.
\end{proof}

\begin{definition}[DimThermalResistance]
  \label{decl:physics:IPhO_2026_4_C_7:DimThermalResistance}
  \lean{IPhO2026Problems.IPhO2026_4_C_7.DimThermalResistance}
  \uses{decl:physics:IPhO_2026_4_C_7:thermalResistanceDimension}
  A unit-independent thermal resistance.
\end{definition}
\begin{proof}
  This is the defining declaration; unfolding it gives the stated typed data or relation.
\end{proof}

\begin{definition}[DimThermalConductivity]
  \label{decl:physics:IPhO_2026_4_C_7:DimThermalConductivity}
  \lean{IPhO2026Problems.IPhO2026_4_C_7.DimThermalConductivity}
  \uses{decl:physics:IPhO_2026_4_C_7:thermalConductivityDimension}
  A unit-independent thermal conductivity.
\end{definition}
\begin{proof}
  This is the defining declaration; unfolding it gives the stated typed data or relation.
\end{proof}

\begin{definition}[DimSpecificHeatCapacity]
  \label{decl:physics:IPhO_2026_4_C_7:DimSpecificHeatCapacity}
  \lean{IPhO2026Problems.IPhO2026_4_C_7.DimSpecificHeatCapacity}
  \uses{decl:physics:IPhO_2026_4_C_7:specificHeatCapacityDimension}
  A unit-independent specific heat capacity.
\end{definition}
\begin{proof}
  This is the defining declaration; unfolding it gives the stated typed data or relation.
\end{proof}

\begin{definition}[DimInverseTime]
  \label{decl:physics:IPhO_2026_4_C_7:DimInverseTime}
  \lean{IPhO2026Problems.IPhO2026_4_C_7.DimInverseTime}
  A unit-independent inverse-time quantity, used for the C.5 graph slope.
\end{definition}
\begin{proof}
  This is the defining declaration; unfolding it gives the stated typed data or relation.
\end{proof}

\begin{definition}[siValue]
  \label{decl:physics:IPhO_2026_4_C_7:siValue}
  \lean{IPhO2026Problems.IPhO2026_4_C_7.siValue}
  The numerical value of a dimension-tagged real quantity in SI units.
\end{definition}
\begin{proof}
  This is the defining declaration; unfolding it gives the stated typed data or relation.
\end{proof}

\begin{definition}[ApparatusGeometry]
  \label{decl:physics:IPhO_2026_4_C_7:ApparatusGeometry}
  \lean{IPhO2026Problems.IPhO2026_4_C_7.ApparatusGeometry}
  \uses{decl:physics:IPhO_2026_4_C_7:DimLength}
  The Figure 17 and procedure geometry.
\end{definition}
\begin{proof}
  This is the defining declaration; unfolding it gives the stated typed data or relation.
\end{proof}

\begin{definition}[Figure17AndProcedureReadout]
  \label{decl:physics:IPhO_2026_4_C_7:Figure17AndProcedureReadout}
  \lean{IPhO2026Problems.IPhO2026_4_C_7.Figure17AndProcedureReadout}
  \uses{decl:physics:IPhO_2026_4_C_7:siValue, decl:physics:IPhO_2026_4_C_7:ApparatusGeometry}
  The geometric and procedural readouts used in C.7.
\end{definition}
\begin{proof}
  This is the defining declaration; unfolding it gives the stated typed data or relation.
\end{proof}

\begin{definition}[PreviousPartC6Data]
  \label{decl:physics:IPhO_2026_4_C_7:PreviousPartC6Data}
  \lean{IPhO2026Problems.IPhO2026_4_C_7.PreviousPartC6Data}
  \uses{decl:physics:IPhO_2026_4_C_7:DimMass, decl:physics:IPhO_2026_4_C_7:DimThermalResistance, decl:physics:IPhO_2026_4_C_7:DimSpecificHeatCapacity, decl:physics:IPhO_2026_4_C_7:DimInverseTime}
  Data reused from the preceding subquestion C.6.
\end{definition}
\begin{proof}
  This is the defining declaration; unfolding it gives the stated typed data or relation.
\end{proof}

\begin{definition}[PreviousPartC6Result]
  \label{decl:physics:IPhO_2026_4_C_7:PreviousPartC6Result}
  \lean{IPhO2026Problems.IPhO2026_4_C_7.PreviousPartC6Result}
  \uses{decl:physics:IPhO_2026_4_C_7:siValue, decl:physics:IPhO_2026_4_C_7:PreviousPartC6Data}
  The reusable C.6 result, stated independently rather than importing a Lean declaration from the previous part.
\end{definition}
\begin{proof}
  This is the defining declaration; unfolding it gives the stated typed data or relation.
\end{proof}

\begin{definition}[ThermalConductionExperiment]
  \label{decl:physics:IPhO_2026_4_C_7:ThermalConductionExperiment}
  \lean{IPhO2026Problems.IPhO2026_4_C_7.ThermalConductionExperiment}
  \uses{decl:physics:IPhO_2026_4_C_7:DimTemperature, decl:physics:IPhO_2026_4_C_7:DimPower, decl:physics:IPhO_2026_4_C_7:DimThermalConductivity, decl:physics:IPhO_2026_4_C_7:ApparatusGeometry, decl:physics:IPhO_2026_4_C_7:PreviousPartC6Data}
  Quantities recorded or modeled in the C-part thermal experiment.
\end{definition}
\begin{proof}
  This is the defining declaration; unfolding it gives the stated typed data or relation.
\end{proof}

\begin{definition}[cylindricalWallAreaMetersSquared]
  \label{decl:physics:IPhO_2026_4_C_7:cylindricalWallAreaMetersSquared}
  \lean{IPhO2026Problems.IPhO2026_4_C_7.cylindricalWallAreaMetersSquared}
  \uses{decl:physics:IPhO_2026_4_C_7:siValue, decl:physics:IPhO_2026_4_C_7:ApparatusGeometry}
  Lateral area “2 π r h” of the active cylindrical wall, in square metres.
\end{definition}
\begin{proof}
  This is the defining declaration; unfolding it gives the stated typed data or relation.
\end{proof}

\begin{definition}[CylindricalConductionLaws]
  \label{decl:physics:IPhO_2026_4_C_7:CylindricalConductionLaws}
  \lean{IPhO2026Problems.IPhO2026_4_C_7.CylindricalConductionLaws}
  \uses{decl:physics:IPhO_2026_4_C_7:siValue, decl:physics:IPhO_2026_4_C_7:ThermalConductionExperiment, decl:physics:IPhO_2026_4_C_7:cylindricalWallAreaMetersSquared}
  The governing physical laws for the experiment.
\end{definition}
\begin{proof}
  This is the defining declaration; unfolding it gives the stated typed data or relation.
\end{proof}

\begin{theorem}[Physics formalization target]
\label{thm:physics:IPhO_2026_4_C_7:target}
\lean{IPhO2026Problems.IPhO2026_4_C_7.acrylicConductivity_from_radial_fourier}
\uses{decl:physics:IPhO_2026_4_C_7:powerDimension, decl:physics:IPhO_2026_4_C_7:thermalResistanceDimension, decl:physics:IPhO_2026_4_C_7:thermalConductivityDimension, decl:physics:IPhO_2026_4_C_7:specificHeatCapacityDimension, decl:physics:IPhO_2026_4_C_7:DimLength, decl:physics:IPhO_2026_4_C_7:DimMass, decl:physics:IPhO_2026_4_C_7:DimTemperature, decl:physics:IPhO_2026_4_C_7:DimPower, decl:physics:IPhO_2026_4_C_7:DimThermalResistance, decl:physics:IPhO_2026_4_C_7:DimThermalConductivity, decl:physics:IPhO_2026_4_C_7:DimSpecificHeatCapacity, decl:physics:IPhO_2026_4_C_7:DimInverseTime, decl:physics:IPhO_2026_4_C_7:siValue, decl:physics:IPhO_2026_4_C_7:ApparatusGeometry, decl:physics:IPhO_2026_4_C_7:Figure17AndProcedureReadout, decl:physics:IPhO_2026_4_C_7:PreviousPartC6Data, decl:physics:IPhO_2026_4_C_7:PreviousPartC6Result, decl:physics:IPhO_2026_4_C_7:ThermalConductionExperiment, decl:physics:IPhO_2026_4_C_7:cylindricalWallAreaMetersSquared, decl:physics:IPhO_2026_4_C_7:CylindricalConductionLaws}
The assigned autoformalize agent should translate this physics problem into Lean declarations in the covered file, with theorem and lemma proof bodies written as `by sorry`.
\end{theorem}
\begin{proof}
This is an autoformalization task, not a proof task. Produce statements that can later be proved by the physics prover without weakening the physical model.
\end{proof}
% --- Archon physics formalization source end ---
